\subsection{Auto-curriculum learning} \label{sec:methods_auto_curriculum}

Given the vastness and diversity of our pre-sampled task pool, it is challenging for an agent to learn effectively with uniform sampling. Most randomly sampled tasks are likely going to be too hard (or too easy) to benefit an agent's learning progress. Instead, we use automatic approaches to select ``interesting'' tasks at the frontier of the agent's capabilities, analogous to the ``zone of proximal development'' in human cognitive development~\citep{vygotsky1978}. We propose extensions to two existing approaches, both of which strongly improve agent performance and sample efficiency (see Section~\ref{sec:results_curriculum}), and lead to an emergent curriculum, selecting tasks with increasing complexity over time.

\paragraph{No-op filtering.} We extend the dynamic task generation method proposed in \citet[Section~5.2]{xland} to our setup. When a new task is sampled from the pool, it is first evaluated to assess whether AdA can learn from it. We evaluate AdA's policy and a ``No-op'' control policy (which takes no action in the environment) for a number of episodes. The task is used for training if and only if the scores of the two policies meet a number of conditions. We expanded the list of conditions from the original no-op filtering and used normalised thresholds to account for different trial durations. See Appendix~\ref{app:autocurriculum} for further details.

\paragraph{Prioritised level replay (PLR).} We modify ``Robust PLR'' (referred to here as \emph{PLR},~\citet{jiang2021robustplr}) to fit our setup. By contrast to no-op filtering, PLR uses a \emph{fitness score} \citep{Schmidhuber1991surprise} that approximates the agent's regret for a given task. We consider several potential estimates for agent regret, ranging from TD errors as used in \citet{jiang2020prioritized}, to novel approaches using dynamics-model errors from AdA (see Appendix~\ref{app:autocurriculum} and Figure~\ref{fig:autocurriculum_abl1}). 

PLR operates by maintaining a fixed-sized archive containing tasks with the highest fitness. We only train AdA on tasks sampled from the archive, which occurs with probability $p$. With probability $1-p$, a new task is randomly sampled and evaluated, and the fitness is compared to the lowest value in the archive. If the new task has higher fitness, it is added to the archive, and the lowest fitness task is dropped. Thus, PLR can also be seen as a form of filtering, using a dynamic criteria (the lowest fitness value of the archive). It differs to no-op filtering in that tasks can be repeatedly sampled from the archive as long as they maintain high fitness. To apply PLR in our heterogeneous task space, we normalise fitness at each trial index by using rolling means and variances, and use the mean per-timestep fitness value rather than the sum, to account for varying trial duration. Finally, since we are interested in tasks at the frontier of an agent's capabilities after across-trial adaptation, we use only the fitness from the last trial. See Appendix~\ref{app:autocurriculum} for further details.



